% ภาคผนวก ก อภิธานศัพท์
\label{app:glossary}

คำศัพท์เฉพาะที่ใช้ในคู่มือและในหน้าจอของระบบมีความหมายดังนี้

\begin{glossarylist}
\term{AI Service}{ส่วนของระบบที่พัฒนาด้วย Python และ FastAPI รับผิดชอบแชตบอต ระบบวางแผนทริป และเครื่องมือ AI ของผู้ดูแลระบบ}
\term{API}{ช่องทางที่โปรแกรมใช้เรียกข้อมูลหรือสั่งงานระบบ ในคู่มือนี้หมายถึง REST API ของ backend}
\term{backend}{ส่วนของระบบที่พัฒนาด้วย Node.js และ Express ให้บริการข้อมูลแก่หน้าเว็บ}
\term{Caddy}{เว็บเซิร์ฟเวอร์ที่ให้บริการหน้าเว็บ ขอใบรับรอง HTTPS อัตโนมัติ และส่งต่อคำขอไปยัง backend และ AI Service}
\term{Cloudinary}{บริการจัดเก็บรูปภาพที่ระบบใช้}
\term{Docker Compose}{เครื่องมือที่สร้างและรันบริการหลายตัวของระบบในคอนเทนเนอร์พร้อมกัน}
\term{Embedding}{เวกเตอร์ตัวเลขที่แทนความหมายของข้อความ ใช้ค้นหาสถานที่ กิจกรรม และฐานความรู้ที่เกี่ยวข้องกับคำถาม}
\term{frontend}{หน้าเว็บที่ผู้ใช้เห็น พัฒนาด้วย React และ Vite}
\term{migration}{ไฟล์ SQL ที่บันทึกการเปลี่ยนแปลงโครงสร้างฐานข้อมูลตามลำดับ}
\term{QR Code}{รหัสสองมิติที่ติดไว้ ณ สถานที่ร่วมรายการ สแกนเพื่อรับพอยท์}
\term{Supabase}{บริการฐานข้อมูล PostgreSQL ที่ระบบใช้}
\term{การซิงก์}{การดึงข้อมูลล่าสุดของสถานที่จาก Google Places มาปรับปรุงข้อมูลในระบบ}
\term{ดัชนีค้นหา}{ชุด Embedding ของข้อมูลทั้งหมดที่แชตบอตใช้ค้นหา ต้องอัปเดตหลังแก้ไขข้อมูล}
\term{ฉบับร่าง}{สถานที่ที่บันทึกไว้แต่ยังซ่อนจากผู้ใช้ รอผู้ดูแลระบบตรวจทานก่อนเผยแพร่}
\term{ฐานความรู้}{เอกสารความรู้ด้านการท่องเที่ยวที่ผู้ดูแลระบบเพิ่มไว้ให้แชตบอตใช้อ้างอิง}
\term{น้องไดโน}{ชื่อของแชตบอตในระบบ}
\term{ผู้ใช้งานทั่วไป}{ผู้ที่เข้าชมเว็บไซต์โดยไม่ได้เข้าสู่ระบบ (Guest)}
\term{พอยท์}{คะแนนสะสมจากการสแกน QR Code ใช้แลกของรางวัล}
\term{ฟิลด์ที่ล็อก}{ข้อมูลของสถานที่ที่ผู้ดูแลระบบแก้ไขเอง ซึ่งการซิงก์จะไม่เขียนทับ}
\term{สมาชิก}{ผู้ใช้ที่สมัครสมาชิกและเข้าสู่ระบบแล้ว (Tourist)}
\term{สถานที่สำคัญ}{สถานที่แนะนำ สถานที่ที่มี QR Code หรือสถานที่ที่มีรายงานค้างอยู่ ซึ่งระบบซิงก์ถี่กว่าสถานที่ทั่วไป}
\end{glossarylist}
